\documentclass[12pt,a4paper]{article}

% \usepackage{array}
% \usepackage{pbox}
% \usepackage{hanging}
\usepackage{fancyhdr}
% \usepackage{caption}
% \usepackage{enumitem}
\usepackage{fancyvrb}
\usepackage{color}
\usepackage[utf8]{inputenc}
\usepackage{graphicx}
\usepackage{listings}
\usepackage{url}

\usepackage{xspace}

\usepackage{rotating}
\usepackage{wrapfig}

\usepackage{cite}
\usepackage{textcomp}

\usepackage[gen]{eurosym}
\usepackage{hyperref}

\usepackage{todonotes}

\usepackage{gensymb}

% \usepackage{draftwatermark}
% \SetWatermarkText{draft v2}
% \SetWatermarkScale{6}

\fancypagestyle{plain}{ %
  \fancyhf{} % remove everything
  \renewcommand{\headrulewidth}{0pt} % remove lines as well
  \renewcommand{\footrulewidth}{0pt}
}

\hyphenation{EISCAT}

\newenvironment{MYitemize}{%
    %% \renewcommand{\labelitemi}{$\rightarrow$}%
    %% \renewcommand{\labelitemii}{$\circ$}%
    %% \renewcommand{\labelitemii}{$\rightarrow$}%
    %% \renewcommand{\labelitemiii}{$\rightarrow$}%
    %% \renewcommand{\labelitemiii}{\colblack $\cdot$}%
    \begin{itemize}}{\end{itemize}}

\newcommand{\bitm}{\begin{MYitemize}}
\newcommand{\eitm}{\end{MYitemize}}

\newcommand{\N}{NeIC\xspace}
\newcommand{\NP}{NeIC~E3DDS\xspace} % NeIC project

\newcommand{\ED}{EISCAT\_3D\xspace}
\newcommand{\EC}{EISCAT\xspace}
\newcommand{\ESA}{EISCAT Scientific Association\xspace}
\newcommand{\SA}{sub-array\xspace}
\newcommand{\SAs}{sub-arrays\xspace}
\newcommand{\RS}{radar site\xspace}
\newcommand{\RSs}{radar sites\xspace}
\newcommand{\OS}{on-site\xspace}
\newcommand{\OC}{operations centre\xspace}
\newcommand{\CC}{control centre\xspace}
\newcommand{\DC}{data centre\xspace}
\newcommand{\DCa}{dCache\xspace}
\newcommand{\DCs}{data centres\xspace}
\newcommand{\UAF}{User Analysis Facility\xspace}

\newcommand{\RB}{ring buffer\xspace}
\newcommand{\fsru}{first stage receive unit\xspace}
\newcommand{\FW}{file writer\xspace}
\newcommand{\SBF}{second beam former\xspace}
\newcommand{\NNB}{$2 \times 100$\xspace} %number of narrow beams

\newcommand{\NBW}{5~MHz\xspace} % narrow bandwidth
\newcommand{\WBW}{30~MHz\xspace} % wide bandwidth

\newcommand{\ramf}{Ramfjordmoen\xspace}
\newcommand{\controllatency}{less than $5$\xspace}
\newcommand{\CP}{C/C++\xspace}
\newcommand{\DI}{DIRAC\xspace}
\newcommand{\DRF}{DigitalRF\xspace}
\newcommand{\HDF}{HDF5\xspace}
\newcommand{\IB}{Infiniband\xspace}
\newcommand{\IBtwo}{200\,Gb/s~\IB}
\newcommand{\Ru}{Rucio\xspace}
\newcommand{\einfra}{e-infrastructure\xspace}
\newcommand{\neinfra}{national e-infrastructure\xspace}
\newcommand{\WR}{White Rabbit\xspace}
\newcommand{\UAP}{User access portal\xspace}

\newcommand{\bfmust}{\xspace{\bf must}\xspace}
\newcommand{\bfshould}{\xspace{\bf should}\xspace}
\newcommand{\Bfmust}{\xspace{\bf Must}\xspace}
\newcommand{\Bfshould}{\xspace{\bf Should}\xspace}
\newcommand{\bfshall}{\xspace{\bf shall}\xspace}
\newcommand{\Bfshall}{\xspace{\bf Shall}\xspace}

\newcommand{\FirstBF}{\xspace{\bf Shall}\xspace}

\newcommand{\visroot}{isualiz}
\newcommand{\vis}{\visroot{ation}\xspace}
\newcommand{\vise}{\visroot{e}\xspace}
% \newcommand{\Vis}{{Visualization}\xspace}
\newcommand{\threed}{{3-dimensional}\xspace}
\newcommand{\GD}{{GUISDAP}\xspace}

\newcommand{\gps}{{Gbit/s}\xspace}

\newcommand{\np}{national provider\xspace}
\newcommand{\nps}{\np{s}\xspace}

\newcommand\EatDot[1]{}

\newcommand{\pc}{prompt computing\xspace}

%\newcommand{\NBW}{5~MHz} % narrow bandwidth
%\newcommand{\WBW}{30~MHz} % wide bandwidth

\newcommand{\nnt}{NeIC NT1\xspace}


% \hyphenation{e-infrastructure}

% \setlength{\topskip}{0mm}
\setlength{\headheight}{15pt}
% \setlength{\topmargin}{-5.4mm}
% \setlength{\textheight}{230mm}
\setlength{\textwidth}{180mm}
\setlength{\oddsidemargin}{-5.0mm}
% \setlength{\evensidemargin}{10.0mm}
% \setlength{\captionmargin}{7mm}

\title{
{\bf Deliverable Document 3} \\
Cluster Management for \ED}
\author{E3DDS Team~\footnote{
Anders Tjulin (EISCAT) {\tt anders.tjulin@eiscat.se};
Ari Lukkarinen (CSC) {\tt ari.lukkarinen@csc.fi};
Assar Westman (EISCAT) {\tt Assar.Westman@eiscat.se};
Carl-Fredrik Enell (EISCAT) {\tt carl-fredrik.enell@eiscat.se};
Dan Johan Jonsson (UiT) {\tt dan.jonsson@uit.no};
Janos Nagy (NSC) {\tt fconagy@nsc.liu.se};
Harri Hellgren (EISCAT) {\tt harri.hellgren@eiscat.se};
Ingemar H\"{a}ggstr\"{o}m (EISCAT) {\tt ingemar.haggstrom@eiscat.se};
Mattias Wadenstein (UmU) {\tt maswan@hpc2n.umu.se};
% Roy Dragseth (UiT) {\tt roy.dragseth@uit.no};
John White (NeIC) {\tt john.white@cern.ch}}}

\date{\today}

\begin{document}

\pagestyle{fancy}
\lhead{\bf E3DDS project}
\rhead{\bf 3: Cluster Management}

\maketitle
\par\noindent
\begin{minipage}{0.5\textwidth}
  \includegraphics[scale=0.18]{NEIC_logo_screen_black.pdf}
  %\vspace{-0.09in}
\end{minipage}
\begin{minipage}{0.5\textwidth}
  \hfill
  %\includegraphics[scale=0.25]{EISCAT3Dlogo1.pdf}
  % New official logo with green text
  \includegraphics[width=0.75\linewidth]{e3d-logo-green-500px}
\end{minipage}

\newpage
\tableofcontents
\newpage

\section{Executive Summary}
\label{exec-summ}

% \todo[inline,color=red]{john to write}
%\todo[inline]{Executive summary to be written at the end of the document writing process? Moved the current exec summary to introduction as that's what it was more like}

This document describes the issues foreseen in deploying and configuring the \einfra for the \ED project. This computing \einfra consists of clusters that will perform specialized tasks.
The management of these clusters should be performed in a robust and audit-able manner using standard tools.
The capacity of the network connections between the sites will influence the deployment of the clusters and therefore their administration.

The possibility to re-configure the clusters from their specialized tasks to perform more generic computing tasks is presented.
This effort should be considered as secondary to the task of running the online data taking.
But some effort should be given to this as the total computing resources, if deployed with a low-latency local area network, would represent a considerable computing power.
This could be used for more difficult analysis tasks such as the 3-D analysis that requires parallel computing capabilities.

Also considered is the means to access these computing clusters and other \ED \newline \einfra.
This takes into account logins ranging from system administrators to public consumers of published data.
Some options are presented that provide services with the correct level of security to allow users to authenticate themselves and login to the appropriate resources.

\section{Purpose} \label{sec:purpose}

The purpose of this document is to describe how the computing clusters for the \ED project can be managed and accessed.
It is intended as a broad overview of the cluster management and access techniques currently used by national \einfra providers.

\subsection{Intended Audience}

This document is intended for the \ED project management and members and also the national computing \einfra providers of the hosting countries.

\section{Introduction} \label{sec:intro}

The \ED sites, the central transmitting/receiving (TX) and remote receiving (RX) sites, will require computing capacity in order to process the data from the First Stage Receive Units (FSRUs) and serve data to the users.
This is expected to be a real-time process as described in E3DDS Project Deliverable 2~\cite{e3dds-del-2} and performed by clusters of computers specialized to the task.


Computing clusters to consider fall into these categories:
\bitm
\item Ring Buffer-Beam Former (RBBF) nodes that house the RAM ring buffer and perform the latter stages of the narrow beam formation;
\item Prompt computing nodes that create higher level time and spatially integrated data products;
\item Cluster containing Monitoring, Radar Control, Databases, Site login;
\item Local site disk buffers for high-speed and short term buffering of low-level data;
\item \ED data archives for long-term storage of data products;
\item User analysis computing nodes.
\eitm
An assumption is made that the TX site, located at Skibotn NO, will function as the ``central'' site as opposed to the two RX sites, located at Kaaresuvanto FI and (nominally) Kaiseniemi SE. % \todo[inline]{CFE: Skibotn and Karesuvanto sites are ready for antenna installation}

The amount of computing required to be deployed at these sites varies according to the capacity and latency of the wide area network (WAN) that connects them.  
The WAN options are described in E3DDS Project Deliverable 1.1~\cite{e3dds-del-1-1} that reports 
% From ~\cite{e3dds-del-1-1}, 
the options are generally a ring topology between the TX and RX sites:
\begin{enumerate}
\item IP-routed connections. 200~Gb/s total capacity;
\item Optical Dense Wavelength Division Multiplexing~\cite{dwdm} (DWDM). 200~Gb/s per RX site to central site;
\item Optical DWDM with ethernet switches on local sites. 4~Tb/s per RX site to central site;
\item Optical DWDM to central site. 6~Tb/s per RX site to central site.
\end{enumerate}

% \subsection{Remote sites} \label{sec:remote}
% reverse the order? 

The deployment of the computing clusters on the RX sites will depend strongly on how all the sites are connected through the Wide-Area network, as described in~\cite{e3dds-del-1-1}.
The computing clusters on the RX sites are expected to be  approximately a subset of the central (TX) site configuration.

In the case of option 1 above, a routed 100~Gb/s WAN between the TX and RX sites, the RX sites must be equipped with identical online computing cluster capabilities (Ring buffer/beam former and prompt computing).
In the other extreme case, option 4, of a 6~Tb/s optical connection between the TX and RX sites, the RX sites do not require any computing capabilities beyond site monitoring and networking switches.

In all cases the management of the above clusters should be designed to be remotely managed, whether from the central TX site or a securely-connected device.


% \subsection{Central site} \label{sec:central}
% \section{IT service management}

% Putting cluster management into a wider context
% Discuss sla, fitsm, itil etc (in particular in relation to external providers) see DO4
% Maybe move this section to DO4?
% fitsm provides sla template for downloading
% MW: there is a tradeoff with a strict process to prevent mistakes/failures vs agile operations that can deploy new stuff with minimal friction. Important to consider the "cost of failure" and chose a solution based on that. Most frameworks have a range of strictness in their implementaiton.

\section{IT service management} \label{sec:itsm}

%define and reference
According to wikipedia, IT service management (ITSM) is the activities that are performed by an organization to design, plan, deliver, operate and control information technology (IT) services offered to customers. 

Some benefits of ITSM:
\bitm
\item Better understand services and responsibilities;
\item Formalize what is expected e.g. through Service Level Agreements (SLAs), Underpinning Agreements~\footnote{An Underpinning Agreement can be seen as a SLA with an external supplier where the service provider is in the customer role.}
and Operational Level Agreements (OLAs);
%define
%An underpinning agreement (UA) can be seen as a service level agreement (SLA) with an external supplier where the service provider is in the customer role.
\item Provision resources in an optimal way (both human and material);
\item Have a common language with external partners;
\item Generally improve professional service delivery.
\eitm
% Implementing ITSM in general also comes with a cost. 
% Human resources are needed for training, implementation, assessment and possibly consultation. 
% Also extra tools are needed to handle issues like, for example, change management and issue tickets.
% configuration database (not only cluster)
%% New text from Mattias
Implementing ITSM in general also comes with a cost.  
Human resources are needed for creating policies, training,
implementation, certification assessment and possibly consultation. 
Also extra tools are needed to handle issues
like, for example, ``Change Management'' and issuing tickets.

A more rigid ITSM reduces the risk of mistakes when making changes to the infrastructure, but increases the amount of work needed to make any changes. 
Risks of not being agile enough to meet needs or frustrated operators departing to a less bureaucratic organisation must be weighed against the cost of failure in various services.

An example of a high ``cost of failure'' is a leak of sensitive information putting the entire organisation and/or its' reputation at risk, compared to an end user computational job getting interrupted and having to re-run their notebook. 
The first event absolutely must not happen, but the second only needs to be rare enough that the scientists are not unhappy with
the service offered.

%Frameworks and tools for ITSM 
In order properly implement ITSM one will need to adopt a standard and/or frameworks for ITSM. 
The most well known is ITIL (Information Technology Infrastructure Library)~\cite{itil},%https://www.axelos.com/best-practice-solutions/itil
the most complete is ISO/IEC 2000~\cite{iso2000} an international standard for IT service management, %https://www.iso.org/standard/70636.html
and the lightweight option is FitSM~\cite{fitsm}.
%https://www.fitsm.eu/
% fitsm standard is free and can be downloaded here https://www.fitsm.eu/downloads/
Strictly speaking neither ITIL or FiTSM are proper standards compared to ISO/IEC 20000, however, both ITIL and FitSM do offer certification.
For processes and process management, this is addressed in the 
% (ITIL 3 ->4, process -> service)
current and latest version of ITIL (version 4)~\cite{itil4}.
% https://www.axelos.com/itil-4

In general for \ED the question will be if or how to implement ITSM.
The ITIL framework is generally aimed towards commercial SLAs where avoidance of legal issues makes the framework heavy for academic purposes.
FitSM is an attempt to be a more lightweight ITSM framework for a broad range of organisations.
In each case, the frameworks certainly contain more than \ED will require so the strategy should be to:
\begin{enumerate}
\item Adopt a framework; 
\item Adapt it to the \ED needs; 
\item Apply the framework as widely as possible.
\end{enumerate}
% \todo[inline]{CFE: Envri-FAIR Task Force 4 is also looking into \url{https://www.coretrustseal.org/} certification of data repositories. Is this applicable / advantageous to us?}

\section{Site Cluster Management} \label{sec:cluster-man}

It will be a complex task to manage and provision the site
clusters. Manually keeping track of hardware configurations and
software stacks with interactive tools is likely to be error prone and
work intensive. The installations should be consistent across the
sites and within each sites there are multiple nodes that ideally
should be configured consistently. We suggest that the complete
configuration and state of the site clusters should be described
through machine readable definition files, which are to be put under
version control, and deployed using dedicated tooling. 
This strategy is more generally known as Infrastructure as code (IaC).

%The configuration is typically stored in a database 
% should also be integrated with ITSM tool


% also comes with a cost similar to FitSm
% also a tradeoff between stricness / agility as commented by MW above

\subsection{Configuration and Management Tools} \label{sec:config}

% make table instead, 
There exist several different frameworks work cluster configuration management. We list here three of the big ones and one particularly geared towards HPC clusters (xCAT):
\bitm
\item xCat~\cite{xcat} open source  Perl, Python, Bash;
\item Ansible~\cite{ansible} enterprise/open source, Python, agentless, INI/YAML configuration, playbook~\cite{ansible-galaxy};
\item Puppet~\cite{puppet} enterprise/open source~\cite{puppet-osp},  C++, Clojure, Ruby Domain Specific Language (DSL) config, modules~\cite{puppet-forge};
\item Chef~\cite{chef} enterprise/open source~\cite{chef-infra}, Ruby, Erlang, DSL config, cookbooks~\cite{chef-cookbooks}.
\eitm
For a more extensive list of configuration management frameworks see~\cite{open-source-config}.
% \todo[inline]{Simon: Something is missing here. JW: Fixed.}
%https://en.wikipedia.org/wiki/Comparison_of_open-source_configuration_management_software
Ansible, Puppet, and Chef have big user communities, where configurations are shared. These shared configurations are referred to as playbooks, modules, and cookbooks respectively. Depending on use case theses can be adopted and/or modified so that one does not need to start from scratch for many basic services. 
In the end, the choice of tool will depend on the choice of operational scenario from E3DDS Project Deliverable DO4~\cite{e3dds-del-4}. It is also possible to use several frameworks together as done at many of the HPC sites in the Nordics. 

%Approaches to managing state: declarative/functional/what vs imperative/procedural/how.
Configuration management is to a large extent about how to manage state. Here, there are two main approaches: \emph{declarative} and \emph{imperative}. In the declarative approach the desired configuration is described and the framework will figure which actions are needed to achieve this. For the imperative approach the configuration explicitly describes the necessary actions to achieve the desired state. In more every day language one could say this is about what state is desired, versus how to achieve that state. 
In programming language terms these approaches closely relate to functional and procedural programming techniques, respectively. 
%Methods to propagate changes: push vs pull, transactions should preferably be idempotent.
There are also two main approaches to propagate changes: \emph{push} versus \emph{pull}. 
Consider server configuration as an example: in the push approach the state will be pushed to the servers from the central configuration data base, whereas in the pull approach the servers will ask for and pull in the state change.
In any case should transactions to propagate state preferably be idempotent, that is, if two identical transaction are issued sequentially the second one should be a null operation.  

Some tools provide roll back to previous known working state, this is usually easier implemented for the declarative approach. 
NixOS~\cite{nixos} is an operating system where this is a built-in feature, but comes with a high learning curve. 

Further, it is desirable that configuration changes are authenticated and auditable. 
To achieve this changes to the configuration database should be signed, and state changes be logged, so that it is possible to later trace and review changes. 

\subsection{Operating System provisioning}
\label{ssec:os-prov}
The overall Operating System (OS) choice for the various clusters should be made based on the technical requirements and also the preference of the operator(s).
Most Linux distributions (Debian, CentOS etc) perform the same tasks with the same level of competence and a high degree of automation.
Other distributions with a lower level of automation e.g. requiring manual compilation and installation of all services (e.g. Arch Linux~\cite{archlinux}),
% \todo[inline]{Carl-Fredrik: Arch uses binary software packages by default, with a build system as an option, but does indeed require manual install and configuration.Gentoo builds software from source. In both cases the builds are automatic though so manual compilation is a misnomer. Manjaro is an interesting option: an Arch derivative with an installer and a repository of frozen packages. Probably best to stick to Ubuntu LTS for security reasons}
can perform tasks far more efficiently but at a cost of access to support.  
There are other distributions aimed at specific types of cluster e.g. the ROCKS~\cite{rocks} for High-Performance computing clusters.
Another consideration is how many other sites, performing the same type of computing, also use the same OS.
This can be important for getting help from others when problems arise.
The choice of OS should weigh these considerations to find the most useful solution that will provide long-term support based on the existing experience of staff.

% \todo[inline,color=red]{JW: How about sections here?}
% \todo[inline,color=red]{JW: Make into sections...}

\subsection{BIOS settings} \label{ssec:bios} % security, vulnerabilities?
%% Unknown?
%% Vendor? ... depends on ... 
The underlying firmware of the clusters, the BIOS (or UEFI), can be updated manually by the administrators.
In the case of a computing cluster with hardware supplied and supported by an external company, user updates (flashing) of the BIOS is not generally accepted practice.
In general, the service contract will stipulate the supplier flashes the BIOS when required.
If remote BIOS flashing is required then the cluster(s) need to be equipped with appropriate hardware e.g. Intel Active Management Technology~\cite{intel-amt}.
The firmware on the First Stage Receive Units can be flashed by \EC staff with the guidance (and supply of the firmware image) from the supplier DA Design~\cite{da-design}.

\subsection{Network configuration} \label{ssec:network}
%% Overall
%% e.g. ROCKS distribution.
%% John some text.
%% DHCP server(s).
%% hardware discovery Xcat or Chef??

The configuration of the network also depends on the choice of WAN connecting the RX and TX sites.
In general the network should be separated into: data network that carries the online data traffic e.g. the UDP packets from FSRU to RBBF nodes and beam files over TCP from RBBF to prompt computing nodes; timing network that carries time signals from the master clock to the timing switches under the subarrays; control network that carries the instructions from the experiment database to the FSRUs and RBBF nodes; monitoring network that carries traffic such as hardware information (temperature, power consumption etc) and video monitoring; back up 4G connection to each RX site.

The network configuration is generally achieved through the creation of sub-networks on the overall private (non externally routable) network over the WAN.
Each entity (ethernet port) must be assigned an IP address within the correct sub-network.
This is usually performed by Dynamic Host Configuration Protocol~\cite{dhcp}  (DHCP).
DHCP can be administered manually, in the case of the UDP data streams from FSRU to RBBF this should be desirable as the FSRU must know exactly which ethernet port to send the correct beam packets.
In the case of the failure of a RBBF ethernet card, the new ethernet MAC address must be entered into the DHCP configuration to receive the correct IP address assignment.

A consideration for network and switch (effectively router nowadays) configuration is that a proprietary solution may provide a Graphical User Interface or some other high-level interface but the configurations may be stored in a proprietary (binary) format.
Subsequent upgrades of software or network hardware can result in loss of configuration data.
A solution that stores configurations in an open, standard format gives resilience to data loss and ease of future upgrades. 
The NAV~\cite{nav} software from Uninett~\cite{uninett} is an option for network and switch configuration.

The back up 4G connection is only used in case of power/network outages.
Therefore the 4G connection only has to route packets to/from the control and monitoring networks.
As it is assumed there is a general power/network outage, data taking and processing cannot proceed.


\subsection{Startup and Shutdown procedures} \label{ssec:power}
%% Dan to write... below...
% this maybe more about cluster operation than cluster configuration

Since the radar sites are remote and possibly not staffed 24/7 it will be extra important to identify which operations will require manual presence at the radar site. In normal daily operation ideally no on-site personal should be necessary. 
%
There are two operation scenarios where it will be difficult to fully describe and automate, even with he help of configuration management tools, that is the startup and shutdown procedures. 

Here, it will be useful to identify which parts of the system depends on which and in which order they should be powered on and off. Probably hands on experience will be necessary to identify all these dependencies both for startup and shutdown. With this in mind, we recommend to attempt a controlled shutdown operation as soon as possible after the site is up and running.

A third operation scenario to consider is service and repair of hardware. One should try to identify in which cases can this be done without a shutdown, a partial shutdown or a if a full shutdown is necessary. 

Finally, one should have a disaster recovery procedure.
Here, all scenarios are almost impossible to foresee, but some obvious ones to consider are power outage, switch failure, node failure etc.

Much of what is described in this section has more to do with cluster and site operation than configuration management. All these scenarios and operations and their possible consequences SLAs and OLAs should be described and documented by the ITSM system.

%% section below moved here from deleted hardware further down

%% Dan to write.
%% keep text general

%% Thinking of good way to not have to go to a node to
%% switch off/on
%% Wake on LAN

%% Lights-out management (LOM) or out-of-band management
% from wikipedia: https://en.wikipedia.org/wiki/Out-of-band_management
%In systems management, out-of-band management involves the use of management interfaces (or serial ports) for managing and networking equipment. 

%% Out-of-band access to servers, switches and if possible FSRUs

%% Many vendors have different implementations (eg HPE iLO)

\subsection{Software provisioning} \label{ssec:software}

The software (radar control, online data processing and user analysis ) that will run on the clusters may be provided through: the OS; a separate software stack; virtual machines or containers; software installation frameworks.
%% \todo[inline]{Carl-Fredrik: Do we have to distinguish system SW, operational SW and user analysis SW? System software installation can be handled with the above tools, Ansible etc, maybe this would also work for RBBF and standard prompt analysis software.}

\subsubsection{CernVM-FS} \label{ssec:cernvmfs}
Software for the RBBF and prompt computing and user computing can be provided by the CernVM File System~\cite{cernvmfs} (CernVM-FS). 
This is a scalable, reliable and low-maintenance software distribution service that is implemented as a POSIX read-only file system in user space. 
It was developed to assist High Energy Physics (HEP) collaborations to deploy software on the worldwide-distributed computing infrastructure used to run data processing applications. 
CernVM-FS transfers data and meta-data on demand and verifies data integrity by cryptographic hashes.
In many cases, it replaces package managers and shared software areas on cluster file systems as means to distribute the software used to process experiment data. 

\subsubsection{EasyBuild} \label{ssec:easybuild}

%% Dan...https://easybuild.readthedocs.io/
EasyBuild~\cite{easybuild} is framework for building and installing software that sits on top of regular building systems like make. 
EasyBuild also generates module files~\cite{lmod}, which make it possible to keep several versions of libraries and software to coexist on the system at the same time. 
The module system achieves this by manipulating the content of the \texttt{PATH} and {\tt{LD\char`_LIBRARY\char`_PATH}} environment variables. %% JW sort this out.
Builds are completely specified by a human readable ``easyconfig'' configuration file which can easily be shared across sites and between users and the official repository can be found at~\cite{easybuild-repo}.

For prompt and user computing to be used in research reproducibility is important. It should be possible for other researchers to repeat the same analysis starting from the same data. Further, it should be possible repeat the same computation, maybe two years later, and still get the same numerical results.
% isolation from base os as far as possible plus exact version specification of dependencies 
EasyBuild also retain logs for traceability of the build processes

A further benefit is that EasyBuild is widely used at HPC centers arround the world, so that the same user computation can readily be performed at a national HPC facility. 
Software built with EasyBuild can also be distributed with the help of CernVM-FS.

There is another framework that is very similar to EasyBuild, called Spack~\cite{spack}, which is also widely used.
% Spack potentially 

\subsubsection{Gentoo project prefix} \label{ssec:gentoo}

The Gentoo Prefix project~\cite{gentoo-prefix} develops and maintains a way of installing Gentoo systems in a non-standard location.
This allows to install Gentoo in another location in the file system hierarchy, hence avoiding conflicts.
By using an offset (the ``prefix'' location), it is possible for many ``alternative'' user groups to benefit from a large part of the packages in the Gentoo Linux Portage tree. 

\iffalse
%% JW we can remove this section
\subsection{Hardware control} \label{ssec:hardware}
%% move to startup and shutdown section
Hardware control: remote console, power management 
%% Dan to write.
%% keep text general

%% Thinking of good way to not have to go to a node to
%% switch off/on
%% Wake on LAN

%% Lights OUt

%% lights out system monitor admin by remote control
%% ILO Integrate lights out system.
%% login to specific management node
%% switch alive

%% Comment out: Power management provided by vendor, server power supplies provide monitoring information connected to the monitoring network.

Remote console
either ssh console to the server or power supply also through the monitoring network.
\fi

% Firmware ... if power supply requires firmware then updated by the vendor ... service agreement.

%% FSRU remote access through network. Control/monitoring.

\subsection{Monitoring and Logging} \label{ssec:logging}

%% Dan

%%Logs vs metrics
%
Performance metrics are typically time-series data, that is, values sampled at regular intervals. 
% metrics examples: failed login attemps last hour, average cpu load, network
%% RRdtool https://oss.oetiker.ch/rrdtool/ RRD file format compact storage of numeric time-series data
%% collectd https://collectd.org/ can collect many types of system level statistics

% logs also for audability
Logs are data with a time stamp to keep track of events, for instance errors, transactions and trace/debug logs. Can vary in format and data volume.

%% John to continue... here. OK?
If logs are to be compared between different services then an important requirement is a standardized logging format. 
At the very minimum, a standard time and server address format will make tasks such as forensics after a security incident possible.
%% Collection

The first step is to decide which and how much data to collect. 
For system events and services {\tt syslog} ({\tt rsyslog}) are standard. 
It is also possible to have other shell scripts and application to log with the logger shell command. 
When it comes to numeric time series data for system resource such as CPU load and memory use, the RRD~\cite{rrdtool} file format offers compact storage for long time series. 
RRDtools~\cite{rrdtool} is a collection of tools to gather, manipulate, and graph rdd files. 
For data collection of system level statistics the {\tt collectd} daemon can also write to the RDD file format.
There are other implementations of Time-Series DataBases that might be used for storing metrics e.g. InfluxDB~\cite{influxdb} and Prometheus' built-in Time-series database~\cite{prometheus}.

% Many other applications

%% Transport to centralized storage
Next, at least some of the logs, should be transported to more centralized storage. 
This can potentially be done in several steps. 
Some data might only be collected and kept on the individual nodes, other data may be stored on a server on site and lastly further data and metrics might be extracted for storage in a centralized data base off the radar sites.
% retries, security
% For logs rsyslog can transport over tcp to a central server. 
{\tt rsyslog} can transport logs over {\tt tcp} to a central server.
Logs can also be kept locally or be temporarily stored in a spool directory before further transport. 
The logging can run asynchronously and supports retries if the central server can not be reached. 
% logstash http://logstash.net/, can aggregate logs from many different sources. has input filter, and output pipilines
In general it might be useful to aggregate logs from many different sources with various formats. One of the most common tools for this purpose is logstash~\cite{logstash}, which comes with a library of filters that can be used in a data processing pipeline. Logstash can also collect and transport logs based on ip addresses. 

%% Storage
% Database
% where
% format, how to access
% For how long time and how much to store

%% Analysis
Finally, logs and metrics should be extracted, analyzed and displayed. One popular tool is the elastic stack~\cite{elastic} based on ElasticSearch for data analysis and Kibana~\cite{kibana} for graphical display. 
Another popular tool to query and vizualise metrics is Grafana~\cite{grafana}. Other popular monitoring tools worth mentioning are Nagios~\cite{nagios}, Zabbix~\cite{zabbix} and Ganglia~\cite{ganglia}.
% https://www.elastic.co/
% Realtime display: Kibana, Grafana https://grafana.com/

%% Alerting
% based on metrics and logging events 
% Grafana
% don't cry wolf

% https://www.nagios.org/
% https://www.zabbix.com/
% http://ganglia.sourceforge.net/
%% Monitoring tools: Nagios, Zabbix, Ganglia

\subsection{Resource Management} \label{ssec:resource}
%% Resource management: Kubernetes? %% Tromso pushes to M$ Azure.
% Use-case... 
% Describe generally Trondheim solution... reference?
% Reboot as slurm daemon nodes... advertise to slurm master?
% could these nodes advertise to somewhere else totally off-site.
% tool to send to different clusters... across admin domains.
% htcondor... pbs, torque, plus... 
% maui
The resources of the clusters, in this case the RBBF and prompt computing nodes, should be managed to provide a benefit to \ED at all times.
When the online data taking is paused, when the radar is off or for other reasons, the clusters should be configured to perform other computing tasks e.g. user analysis of data.
These nodes can be configured to run the online data handling on ``bare metal"~\footnote{From wikipedia: In computer science, bare machine (or bare metal) refers to a computer executing instructions directly on logic hardware without an intervening operating system.} and, when the online data is paused, run other computing tasks within other environments.

A simple configuration is to start a slurm~\cite{slurm} daemon, for a cluster management and job scheduling system, that advertises to a central job queue (slurm master) and accepts computing jobs.
In a similar manner there are other job submission systems that can be employed e.g. Portable Batch System~\cite{pbs}, TORQUE~\cite{torque} or HTCondor~\cite{htcondor}.
An attractive feature of this approach is the relative ease of deploying multi-node e.g. using MPI~\cite{mpi} computing jobs.

Another approach can use virtualized environments.
When the online data-taking is paused, the cluster nodes launch virtual machine(s) or containers that have a standard \ED software and configuration environment.
This can be achieved through various levels of virtualization ranging from the low-level xen~\cite{xen}, kvm~\cite{kvm} or QEMU~\cite{qemu} to mid-level libvirt~\cite{libvirt} to higher-level OpenStack~\cite{openstack} or Kubernetes~\cite{kubernetes}
(running Docker~\cite{docker} or Singularity~\cite{singularity}) containers.
%% Kubernetes on top.
The lower level solutions require more expertise in low-level operating systems issues but can be operated with lower \einfra and personnel resources.
Whereas the higher level solutions hide the underlying complexity but require hard-to-find expertise to install and maintain the installation.
% John to describe libvirt possibility. slurm daemon etc. i.e. not openstack 
% etc.
Some considerations for resource management is the lengths of time that the online data taking is to be paused and the expected warning time to re-start the online data taking.

If the data taking pauses have durations of approximately the same length or longer than a user analysis job then jobs can be completed. 
Whereas if the data taking pauses are much shorter than the users' analysis jobs then these will need to be paused, stored and then re-started on the next pause.
This adds time to stop or pause a job.
If needed, jobs can be cancelled in a batch system with the loss of the results to that point.

Virtual machines and containers can be started/stopped in the order of seconds~\footnote{Containers can start/stop reliably in the order of sub-seconds.} with the virtual machine state preserved.
If the warning time to restart online data taking is shorter than e.g. 30 seconds then frequent updates of the virtual machine (or container) state to storage will be required in order to facilitate a fast shutdown and return to bare-metal online data processing.

At the time of writing this document, the frequency and length of the data taking pauses and the expected average length (or distribution) of user jobs are not known.
Also, similarly the warning time required to re-start data taking is not known.

% Dan also write caveat about resources?
Finally, we like to note that using RBBF and prompt computing nodes for user compute, also comes with a cost in terms of human resources in order to setup and maintain this alternate use case.  
%Not strictly needed for the operation of the radar. Get system up and working first then consider resource reuse.

\section{Cluster Login}
\label{sec:cluster}

The data processing site (initially most likely the Skibotn site) will host several distinct clusters for \ED.
The Ring Buffer-Beam Former (RBBF) nodes will function as the RAM ring buffer and also produce the narrow angle beams.
These nodes will also write the narrow angle beams into files to be forwarded to the prompt computing for further processing into data products.
The RBBF cluster is expected to be composed of $30$~\footnote{At the time of writing, it is expected that each CPU can handle one wide-angle beam from the FSRUs. Therefore 10 dual socket servers are required for each RX site.} high-RAM and throughput computing nodes~\cite{amd-epyc}.

The ``Prompt Computing'' is expected to consist of $N$ commodity computing nodes (optimized for reprocessing of files) and/or $M$ (where $N > M$) nodes containing GPU accelerators.

It is also expected that another system, used for administrative purposes, will be sited at Skibotn. This administrative cluster is not expected to be anywhere near the same scale as the RBBF or Prompt computing clusters
and is a good candidate to be hosted within virtual machines.
The RBBF and prompt computing clusters are planned to be used for \ED user analyses during the times when the systems would otherwise be (partially) idle.

If the Optical DWDM fibre ring design is selected, the main \ED compute resources will only be accessible through this network. Routing between \ED and the public internet will be managed by a single access router and firewall, located at the Skibotn site or elsewhere along the ring.
% Considering also e.g. a location close to Eiscat HQ in Kiruna
This will be the only way to access the clusters and will allow login (e.g. ssh) and file transfer (FTS, etc) only from other \EC sites. 

\EC normal users can access non-sensitive data under the \EC data embargo rules, and submit requests to run experiments. 
These users should never log in to the online site clusters.
Given this, there are several types of logins to consider with different access mechanisms:
\bitm
\item System administrator: login with full privileges for system maintenance. Typically these are root privileges that can affect anything on the system.
% This would typically be root

\item Radar developers: login with sufficient capabilities to install and modify radar software.
% comparable to "kstdev" on legacy systems?

\item Radar expert operators: radar users and developers with non-root privileges to control all aspects of the transmitters and receivers in realtime, but not to change the operating system or install software.
% "eiscat"  on the legacy systems

\item Machine users: Other systems given privileges to change the scheduling of predefined radar modes in real time. 

\item \EC privileged data users logging in to access sensitive lower-levels of data.

\item Data manager machine login: Sufficient privileges to read and replicate non-sensitive data to external sites and perform machine to machine logins.

% \todo[inline,color=red]{JW: Surely data replication will be automatic? Given the volumes of data expected and the expertise existing to set up such an automatic system? Could this be re-written as ``Data management machine login"??}

\item Analysis user: Sufficient privileges to start VMs and stage data on connection from the user analysis job submission portal.
 
\eitm

\subsection{System Administrator Login}
\label{sec:admin}

% Here describe the NT1 solution for admin logins...
% 2 factor
% other candidates?

The administrators of these clusters will need to be able to access the \ED systems through the Optical DWDM network via the access router.
This is different to the logins described in Section~\ref{sec:users} as these will be performed by \EC employees or contractors
in order to repair/configure/update the systems.
It is important to remember that these logins will be requesting privileged access (up to root level on Unix-type systems) and contrast to user logins on other \einfra resources.
These administrators may be logging in to the clusters through local terminals or remotely via the internet from any \EC member institution
or contracted partner.
The second case is far more likely.
Some basic requirements for administrative logins can be given:
\bitm
\item Login over secure connections only. i.e. no unencrypted data anywhere

\item Login restricted to known set of users; highest level of authorization required, e.g. by two-factor authentication~\cite{two-factor}

\item Logins and actions traceable (logged)
\eitm

% \subsubsection{Bastion hosts} \label{ssec:bastion}
\subsubsection{Signed ssh credentials via bastion host} \label{ssec:bastion}

This subject of securely logging in multiple administrators over many nodes has been studied in industry~\cite{secure-ssh}.
A solution~\cite{fb-ssh} has been implemented by industry leader Facebook using open-source tools.
This approach has been implemented in \nnt, an academic computing \einfra, that shares the same characteristics as the future \ED \einfra.

% \todo[inline]{FIX}
This scheme is based on a key-pair-enabled OpenSSH server, an implementation of the industry-standard Secure Shell (SSH) for accessing systems.
This scheme takes advantage of its special features, that are not widely leveraged, to provide both security and reliability using signed certificates with principals.

This solution is discussed in detail in~\cite{fb-ssh}.
The \nnt deployment~\cite{nt1-ssh} also uses the ``Bastion host" concept where the administrator logs in to a highly-protected secured host.
Engineers (administrators) authenticate to the bastion host(s) and obtain the signed SSH certificate.
This certificate contains all principals (information) allowed for the specific administrator. 
Two-factor authentication can be implemented on the Bastion host to ensure an improved level of strong authentication.
This certificate can now be used to SSH into cluster nodes as either root or other lower privileged users.
In this fashion, no administrator has more access than they require~\cite{privileges}.

\subsubsection{Local login node}
\label{ssec:local-login-node}

Another scheme for administrators is to deploy (bare metal or virtual) a node, network-local to the \ED clusters, to act as a login node.
The administrator logs in to the login node using SSH.
Once logged in to this node, they can then further login to the cluster(s) via a second SSH step.
The login node can also be enabled to require two-factor authentication.

\subsubsection{Direct logins}
\label{ssec:direct}
% Carl-Fredrik: Isn't this equivalent to opening the firewall between the public internet and the DWDM ring to the world? That should be avoided.
Another option is for system administrators to directly log in to the cluster nodes.
This would require the nodes to be open to the internet and the need to maintain a login directory (LDAP or other) to synchronize user name and passwords between nodes.
Root privileges would be acquired through privilege escalation (sudo or other) after an administrator has logged in.
This scheme can be considered to be the least secure.

\subsection{Radar Developer and Expert Operator Login}
\label{sec:operator}
The operators of the radar system need access to the scheduler, FSRUs, RBBF and prompt computing nodes in order to update the parameters of the system.  The privileges required for these operations should not be at the highest level (root), to avoid operations that could affect the basic functionality of the systems.

The radar operators should have enough privileges on the clusters to be able to change parameters in databases to affect the beam-forming and change the common program operations in the prompt computing cluster.
Developers must also be able to update or change the software.
The interactions of the radar operators  with the system are expected to be a combination of graphical (Jupyter notebooks or other) and command-line (e.g. ssh).

In the case of the radar operators' graphical interaction with the system, step-up authorization with strong authentication will be required, discussed in Section~\ref{sec:users}.
For command-line interactions, the radar operators should login using the same method as the system administrators in Section~\ref{sec:admin}.
The radar operator would then receive only the correct level of privilege to do the task.

\subsection{Machine Logins}

Other computer systems within the \ED \einfra will have to access the site computing and storage for at least these purposes:
\bitm
\item Data replication and staging (DIRAC, Rucio, etc TBD);
% \todo[inline,color=red]{JW: Surely data replication will be automatic? Given the volumes of data expected and the expertise existing to set up such an automatic system?}
\item Realtime rescheduling of experiment modes through the scheduler API.
\eitm

The authorization mechanism here is likely to be X.509 proxy certificates with a limited lifetime.
These credentials, robot certificates, should be automatically generated to be used by the processes.
% Add info about connections between NT1 systems etc?
This should never involve direct connections to the RBBF or prompt compute clusters, only to site storage and to the remote scheduler API, which should run on the control system.
% or even on an external gateway host outside the e3d network?



\subsection{User Login}
\label{sec:users}

The NeIC E3DS project has investigated~\cite{e3ds-md3}
the requirements for users functioning in a scientific (as opposed to administrative) role.
These users need to gain access to the various levels of \ED data ranging from raw level~1 to physical parameters level~3 data.
The \ED data levels are defined in Table~1 of E3DS Project Deliverable MD-3~\cite{e3ds-md3}.
Users need to authenticate (identify) themselves in order to access the data and \einfra resources.
The degree that an authentication can be trusted is given by its strength~\cite{owasp-aai} ranging from weak (e.g. user name and password) to strong (e.g. two-factor authentication from independent sources).

Depending on the data level requested to be accessed by the \ED user, the strength of the authentication may vary.
The level~1 (a or b) data may be more closely controlled as it may % contain hard-target information e.g. satellites or other objects that authorities may have reservations about.
possibly contain more sensitive information subject to regulation.
The higher data levels (to level 3 or 4) are integrated summary or user-produced data and may not require such stringent authentication.

\subsubsection{Normal Users}
\label{ssec:normal}

As outlined in E3DS Project Deliverable MD-3~\cite{e3ds-md3}, the expected primary access to \ED data by users is expected to be through a data portal e.g. DIRAC~\cite{dirac}.

The standard method to authenticate to the DIRAC portal is by possession of a recognized X.509 certificate, either installed to a browser or presented through a command line interface.
This authentication method entails an operational overhead~\footnote{For a small entity like \EC expecting an order of magnitude more users than currently, this will probably be a heavy overhead.} in certificate issuance and maintenance.
In order to use this authentication scheme, the \ED users will also have to be educated in the usage and possession of X.509 credentials.
A hard task~\cite{cert-example} for any research community that mainly consists of scientists rather than \einfra experts.

A more lightweight and familiar (for the end users) authentication method is available.
The ``Federated Identity" concept uses the concept of a ``trusted third party" that can provide authentication credentials.
In this scheme trusted institutional Identity Providers (IDPs) are given permissions (federated) to provide authentication credentials to each other. 
An example of such a scheme is the EduGAIN~\cite{edugain} federation.

These credentials are passed from the institutional IDPs to services 
through an {\bf A}uthentication and {\bf A}uthorisation {\bf I}nfrastructure (AAI).
Typically a user logs in to their home institute IDP with their home institute user name and password (or similar) or even using commercial social media accounts.
The IDP generates the necessary identity tokens (X.509 certificates or other).  
These tokens are then transferred transparently to e.g. the DIRAC portal and downstream services that the portal contacts or uses.
Therefore the users in a authentication-federated institute may access services~\footnote{The authorization permissions and actions for a remote user presenting federated authentication credentials rests with the local host.} in another federated institution.

The DIRAC portal is an example of a service that integrate with proxy authentication services such as EGI Checkin~\cite{egi-checkin} or eduTEAMS~\cite{eduteams}.
These services can provide the full AAI by managing group membership information, managing groups and subgroups, add and remove users, manage user consent and the acceptable usage policy.
Both the EGI Checkin~\cite{egi-checkin} and GEANT eduTEAMS~\cite{eduteams} have plans to incorporate the multi-factor authentication discussed below in Section~\ref{ssec:high}.

\subsubsection{Privileged Users}
\label{ssec:high}

Some \EC users may be considered to have special privileges. 
There will be users who are allowed to request (or even start) an experiment, and authorized users must also be able to upload analysis software for running in an isolated VM environment. These operations are expected to take place through Jupyter~\cite{jupyter} notebooks.

Another important consideration for user privilege to access the \ED data are requirements that come from external sources e.g. funding and other governmental agencies.
If the identity of users who access the sensitive (level~1) \ED data is required to be known to a higher strength than that provided by an academic identity federation, then other steps can be taken.
An AAI that provides a higher strength of authentication, typically using a multifactor scheme, can be used.

Multifactor authentication works with two separate security or validation mechanisms. 
Typically, one is a physical validation token, and one is a logical code or password and both must be validated before accessing a secured service or product.
The user first authenticates to their home IDP, typically by password-based authentication (vulnerable to various attacks), and subsequently by a second authentication factor required by the AAI.
This can be a SMS One-Time Password~\cite{otp}, RSA token~\cite{rsa-token} or a mobile app etc.

The ELIXIR AAI~\cite{elixir-aai} is an example of an AAI that enables researchers to use their home organisation credentials or community or commercial identities (e.g. ORCID, LinkedIn) to sign in and access data and services.
Importantly, ELIXIR AAI includes a ``step-up" authentication level for services with sensitive data, that expect not just an authentication credential but also authentication which relies on multiple authentication factors.

Usage of the ELIXIR AAI by \ED may be problematic as it looks likely that there may be policy and governance problems.
ELIXIR AAI is designed for access to resources/services for professional activities related to the development and use of services for life science research.
Therefore use of EGI Checkin or eduTEAMS with multifactor authentication enabled would solve the authentication issues for \ED.

\subsubsection{Public Users} \label{ssec:public}

Another class of users can access publicly available datasets.
These are typically un-embargoed high-level data sets or data sets made public via publications.
In this case, the users can be any member of the general public.
A high level of authentication is not required to access these data sets but access should not be completely anonymous.
A capability to identify unique user accesses to data sets is desirable for accounting and reporting to funding agencies etc.
Therefore, the selected user authentication system should be able to authenticate using ``lowest-common denominator'' credentials e.g. social media or online email accounts etc.

% Some one who requests an experiment.
% Users request experiment
% 


% end of the text of the document...

\newpage
\bibliography{main}{}
\bibliographystyle{unsrt}

\end{document}
